\documentclass[12pt,a4paper]{exam}

% --------------------
% Packages
% --------------------
\usepackage[utf8]{inputenc}
\usepackage[T1]{fontenc}
\usepackage{bm}
\usepackage[french]{babel}
\usepackage{amsmath,amssymb}
\usepackage{geometry}
\usepackage{multicol}

\usepackage{pgfplots}

\pgfplotsset{compat=1.18}
\geometry{margin=2cm}

% --------------------
% En-tête / pied (exam)
% --------------------
\header{}{Contrôle de mathématiques}{4\up{ème}1}
\footer{}{\thepage}{\it{G.M.}}

% --------------------
\begin{document}

\begin{center}
    \Large\textbf{Devoir surveillé\\Nombres premiers, puissances et fractions}
\end{center}

\vspace{0.3cm}

\noindent
\textbf{Durée :} 1h \hfill
\textbf{Calculatrice :} \underline{NON} autorisée

\vspace{20pt}

% ======================================================
% =================== EXERCICE 1 =======================
% ======================================================
\noindent
\textbf{Exercice 1 : Nombres premiers et calculs (\ \ \ \ \ \  / 6 points)}
\vspace{5pt}

\begin{questions}

    \question[1]Donner la \underline{définition} d’un nombre premier.
    
    \question[1]Donner quatre nombres premiers \textbf{supérieurs à $\bm{11$}}.

    \question[4]Décomposer les nombres suivants en produit de facteurs premiers :
    \begin{parts}
        \part $90$
        \part $840$
    \end{parts}

\end{questions}

\vspace{20pt}


% ======================================================
% =================== EXERCICE 2 =======================
% ======================================================
\noindent
\textbf{Exercice 2 : Puissances (\ \ \ \ \ \  / 9 points)}
\vspace{5pt}

\begin{questions}

    \question[3]Effectuer les calculs suivants :
    \begin{multicols}{2}
    \begin{parts}
        \part $(-8)^2 + 3^2$
        \part $4^3 - 10^0$
    \end{parts}
    \end{multicols}
    \begin{multicols}{2}
    \begin{parts}
        \setcounter{partno}{2}
        \part $-5^2 \times (4 - 3)^2$
    \end{parts}
    \end{multicols}

    \question[3]Écrire sous \underline{\textbf{forme décimale}} les calculs suivants :
    \begin{multicols}{2}
    \begin{parts}
        \part $3,68 \times 10^4$
        \part $1333,4 \times 10^{-6}$
    \end{parts}
    \end{multicols}
    \begin{multicols}{2}
    \begin{parts}
        \setcounter{partno}{2}
        \part $(-5)^2 \times 10^{-4}$
    \end{parts}
    \end{multicols}

    \question[3]Écrire les nombres suivants en \underline{\textbf{notation scientifique}} :
    \begin{multicols}{2}
        \begin{parts}
            \part 35 000
            \part 0,000873
        \end{parts}
    \end{multicols}
    \begin{multicols}{2}
        \begin{parts}
            \setcounter{partno}{2}
            \part 4 220 000
        \end{parts}
    \end{multicols}
    
\end{questions}

\vspace{20pt}

% ======================================================
% =================== EXERCICE 3 =======================
% ======================================================
\noindent
\textbf{Exercice 3 : Fractions (\ \ \ \ \ \  / 5 points)}
\vspace{5pt}

\begin{questions}

    \question[2] Effectuer les calculs suivant et simplifier si possible
    \begin{multicols}{2}
        \begin{parts}
            \part $\dfrac{5}{11} + \dfrac{3}{-4}$
            \part $1 - \dfrac{6}{49} - \dfrac{4}{7}$
        \end{parts}
    \end{multicols}

    \vspace{5pt}
    
    \question[3] Adrien affirme qu'il a parcouru les $\dfrac{7}{9}$ du parcours accro-branche lors 
    de la sortie scolaire. Léa, elle, dit qu'elle en a fait les $\dfrac{4}{5}$. Elle prétend l'avoir battu. 
    \vspace{5pt}
    
    A-t-elle raison ?

\end{questions}

\vspace{20pt}

% ======================================================
% =================== EXERCICE 4 =======================
% ======================================================
\noindent
\textbf{Exercice 4 (facultatif) : \underline{BONUS} (\ \ \ \ \ \  / 2 points)}
\vspace{5pt}

Le coeur humain effectue environ $5\ 000$ battements par heure.

\begin{questions}

    \question Écrire $5\ 000$ en notation scientifique.
    \question Calculer le nombre de battements effectués en un jour, sachant qu'une journée 
    dure 24h.
    \question Calculer le nombre de battements effectués pendant une vie de 70 ans. On considère 
    qu'une année dure 365 jours.
    \vspace{5pt}

    Donner ce résultat en notation scientifique.

\end{questions}

\vspace{20pt}

\begin{center}
    \textit{Fin du devoir}
\end{center}

\end{document}
